During the current reporting period at Universitat Jaume I, my doctoral research has focused on context-aware ubiquitous computing and digital mental health interventions. While current Ecological Momentary Interventions and Just-in-Time Adaptive Interventions rely heavily on internal physiological tracking and self-reported symptoms, my work investigates how external contextual factors can be leveraged to deliver more situation-aware psychological support.
The primary milestone achieved during this period is the execution and drafting of a comprehensive systematic review titled "The Role of Context-aware Technological Devices and External Context in Mental Health Interventions: A Systematic Review". Following the PRISMA 2020 guidelines, I conducted a structured search and screening workflow across four major scientific databases (Scopus, Web of Science Core Collection, APA PsycInfo, and PubMed), with eligibility reviews synthesizing 40 peer-reviewed studies.
Through this synthesis, I developed a multidimensional classification framework analyzing targeted mental disorders, consumer-grade hardware platforms, external context types (physical, social, and computing/temporal), capture techniques, clinical purposes, delivery mechanisms, and evaluation designs.
The findings show a strong predominance of passive computing and temporal telemetry, driven by mobile operating system constraints and privacy considerations, as well as a clear reliance on user-driven and rule-based delivery mechanisms over black-box automated algorithms.
Additionally, the review identifies critical translational gaps, including the scarcity of multi-modal external sensing and the decline of usability evaluations in mature clinical trials.
The manuscript is currently being revised, with narrative refinements and reproducibility archiving underway before peer-reviewed journal submission. Building directly on the empirical gaps established in this review, the upcoming phase of my doctoral research will focus on designing, implementing, and evaluating a situation-aware mobile intervention prototype that integrates passive external context telemetry with interpretable, privacy-preserving therapeutic delivery.